% ==========================================================================
%  Chapter 53 — Corotational Dynamic RC Building: Formulation,
%               VTK Reconstruction Fix, and Kinematic Incompatibility
% ==========================================================================

\chapter{Corotational Dynamic RC Building Example}
\label{ch:corotational-dynamic-building}

This chapter documents the transition of the five-story RC building
example from small-rotation static analysis
(Chapter~\ref{ch:regular-5story-rc-nonlinear-building}) to a
\emph{corotational dynamic} formulation.  Three major topics are covered:
\begin{enumerate}
  \item the corotational beam kinematics and their integration into the
        existing structural pipeline;
  \item a critical bug found and fixed in the VTK displacement-field
        reconstruction for corotational elements;
  \item an analysis of the kinematic incompatibility between linear
        Mindlin--Reissner shells and corotational beams, and its
        consequences for numerical stability.
\end{enumerate}


% ══════════════════════════════════════════════════════════════════════════
\section{Motivation and Scope}
\label{sec:corot-motivation}

The small-rotation formulation used in
Chapter~\ref{ch:regular-5story-rc-nonlinear-building} treats each beam
element as geometrically linear: the transformation matrix~$\mathbf{T}$
is evaluated once from the reference configuration and is never updated.
This is adequate for load levels where each storey drift remains well
below $1\%$.  For seismic or extreme-wind scenarios, however, the
building can develop inter-storey drifts of $2$--$5\%$ or more, at which
point the geometric nonlinearity encoded in $P$-$\Delta$ effects and
large rotations becomes essential.

The current chapter therefore upgrades the frame elements to the
\texttt{beam::Corotational} kinematic policy while keeping the same
material model (nonlinear fiber sections), the same global topology, and
the same \texttt{DynamicAnalysis} PETSc time-stepping infrastructure.
The floor slabs remain elastic Mindlin--Reissner shells.


% ══════════════════════════════════════════════════════════════════════════
\section{Corotational Beam Kinematics}
\label{sec:corot-kinematics}

% ------------------------------------------------------------------
\subsection{Conceptual Framework}

The corotational formulation decomposes the total motion of a beam
element into:
\begin{enumerate}\itemsep0pt
  \item a \emph{rigid-body} component (translation and rotation of the
        element chord), and
  \item a \emph{deformational} component (local strains measured in the
        corotated frame).
\end{enumerate}

The fundamental idea is that a small-strain $B$-matrix evaluated in the
\emph{rotating} local frame captures effectively the same physics as a
geometrically exact formulation, at a fraction of the implementation
complexity.

Let $\mathbf{x}_0, \mathbf{x}_1 \in \mathbb{R}^3$ be the reference
nodal positions and $\mathbf{u}_0, \mathbf{u}_1$ the corresponding
global displacement vectors.  The \emph{current} nodal positions are
\begin{equation}\label{eq:current-positions}
  \hat{\mathbf{x}}_i = \mathbf{x}_i + \mathbf{u}_i,
  \qquad i = 0, 1.
\end{equation}

The current chord vector is
$\hat{\mathbf{d}} = \hat{\mathbf{x}}_1 - \hat{\mathbf{x}}_0$, with
current element length
\begin{equation}
  l_n = \lVert \hat{\mathbf{d}} \rVert.
\end{equation}

% ------------------------------------------------------------------
\subsection{Corotated Frame Update}

In three dimensions the corotated frame $\mathbf{R}_n$ is constructed as
follows:
\begin{align}
  \mathbf{e}_1 &= \frac{\hat{\mathbf{d}}}{l_n},
  \\[6pt]
  \tilde{\mathbf{e}}_2 &= \mathbf{e}_{2,\mathrm{ref}}
    - \bigl(\mathbf{e}_{2,\mathrm{ref}} \cdot \mathbf{e}_1\bigr)\,
      \mathbf{e}_1,
  \\[4pt]
  \mathbf{e}_2 &= \frac{\tilde{\mathbf{e}}_2}
                       {\lVert \tilde{\mathbf{e}}_2 \rVert},
  \\[4pt]
  \mathbf{e}_3 &= \mathbf{e}_1 \times \mathbf{e}_2.
\end{align}

Here $\mathbf{e}_{2,\mathrm{ref}} = \mathbf{R}_0^{\mathsf{T}} \mathbf{i}_2$ is
the reference cross-section $y$-axis.  The projection onto the plane
normal to the current chord ensures smooth frame evolution even for
moderate rotations.  In the degenerate case
$\lVert\tilde{\mathbf{e}}_2\rVert < 10^{-12}$ the algorithm falls back to
the reference-vector construction.

In the implementation, frame update is \emph{embedded} inside every
element method that requires the current configuration:
\begin{lstlisting}
// BeamElement.hh
Eigen::VectorXd compute_internal_force_vector(const Eigen::VectorXd& u_e)
{
    if constexpr (!KinematicPolicy::is_geometrically_linear) {
        update_current_frame(u_e);   // recompute R_n from (x+u)
    }
    auto T = transformation_matrix();    // T = blkdiag(R_n, R_n)
    auto u_loc = KinematicPolicy::extract_local_dofs(u_g, frame0_, frame_n_, T);
    // ... B_local, sigma, f_loc ...
    return (T.transpose() * f_loc).eval();
}
\end{lstlisting}
This design ensures the frame is self-consistently updated at every Newton
iteration without requiring explicit calls from the analysis driver.

% ------------------------------------------------------------------
\subsection{Deformational DOF Extraction}

The \texttt{beam::Corotational::extract\_local\_dofs} function strips the
rigid-body motion from the global DOFs.  For a 3D two-node beam with 12
global DOFs $[\mathbf{u}_1,\bm{\theta}_1,\mathbf{u}_2,\bm{\theta}_2]$:

\begin{equation}\label{eq:deformational-dofs}
  \mathbf{u}_\mathrm{def} =
  \begin{bmatrix}
    0 \\ 0 \\ 0 \\
    \bm{\theta}_{1,\mathrm{loc}} - \bm{\alpha} \\
    \Delta l \\ 0 \\ 0 \\
    \bm{\theta}_{2,\mathrm{loc}} - \bm{\alpha}
  \end{bmatrix},
\end{equation}
where $\Delta l = l_n - l_0$ is the axial elongation, and
$\bm{\alpha}$ is the rigid-body rotation vector extracted from the
relative rotation between $\mathbf{R}_n$ and
$\mathbf{R}_0$.

Crucially, node~1 has \emph{zero} translational deformation (it is the
local origin) and node~2 has only an axial component $\Delta l$.  This
has fundamental consequences for post-processing, as discussed in
\cref{sec:vtk-displacement-fix}.

% ------------------------------------------------------------------
\subsection{Tangent Stiffness and Geometric Contribution}

The tangent stiffness in global coordinates is
\begin{equation}\label{eq:corot-tangent}
  \mathbf{K}_t
  = \mathbf{T}^{\mathsf{T}} \underbrace{\sum_{g=1}^{n_g}
      w_g \, \lvert J \rvert_g \;
      \mathbf{B}_g^{\mathsf{T}} \,
      \mathbf{C}_{t,g} \,
      \mathbf{B}_g}_{\mathbf{K}_\mathrm{loc}}
    \mathbf{T}
  +
  \mathbf{T}^{\mathsf{T}} \, \mathbf{K}_\sigma \, \mathbf{T},
\end{equation}
where $\mathbf{K}_\sigma$ is the geometric stiffness due to the element
resultants.  Both terms are computed in
\texttt{BeamElement::compute\_tangent\_stiffness\_matrix}; the
geometric stiffness is enabled by the compile-time flag
\texttt{KinematicPolicy::needs\_geometric\_stiffness = true} for the
corotational policy.


% ══════════════════════════════════════════════════════════════════════════
\section{Dynamic Analysis Configuration}
\label{sec:corot-dynamic-config}

The dynamic analysis uses the same \texttt{DynamicAnalysis} class
documented in Chapter~\ref{ch:phase5-dynamic-seismic}, with
PETSc's \texttt{TSALPHA2} generalized-$\alpha$ integrator.

% ------------------------------------------------------------------
\subsection{Loading and Ramps}

To avoid dynamic shock and ensure a gradual transition to the loaded
state, the external force function is defined as
\begin{equation}\label{eq:force-function}
  \mathbf{f}_\mathrm{ext}(t)
  = g(t)\,\mathbf{f}_\mathrm{grav}
  + h(t)\,\sin\!\bigl(\tfrac{2\pi t}{T_\mathrm{exc}}\bigr)\,
    \mathbf{f}_\mathrm{lat},
\end{equation}
with the ramp functions
\begin{equation}
  g(t) = \min\!\bigl(\tfrac{t}{t_\mathrm{grav}},\; 1\bigr),
  \qquad
  h(t) = \begin{cases}
    0 & \text{if } t < t_\mathrm{grav}, \\
    \min\!\Bigl(\dfrac{t - t_\mathrm{grav}}{t_\mathrm{lat} - t_\mathrm{grav}},\; 1\Bigr) & \text{otherwise.}
  \end{cases}
\end{equation}

The gravity ramp $g(t)$ linearly increases over $t_\mathrm{grav} = 1.0$\,s
to avoid a step function in the mass-proportional gravity load.  The
lateral cyclic loading $h(t)$ is delayed until the gravity ramp
completes and then smoothly increases over the interval
$[t_\mathrm{grav},\; t_\mathrm{lat}]$ with $t_\mathrm{lat} = 1.5$\,s.

% ------------------------------------------------------------------
\subsection{Time-Integration Parameters}

\begin{table}[htbp]
\centering
\caption{Dynamic analysis parameters for the corotational RC building.}
\label{tab:corot-params}
\begin{tabular}{lrl}
\toprule
  \textbf{Parameter} & \textbf{Value} & \textbf{Notes} \\
\midrule
  Spectral radius $\rho_\infty$ & 0.9 & High-frequency dissipation \\
  Time step $\Delta t$ & 0.005\,s & 1000 steps for $T_f = 5\,$s \\
  Adaptive stepping & off & \texttt{TSADAPTNONE} \\
  Max SNES failures & 5 & Allow temporary divergence \\
  KSP type & \texttt{preonly} & Direct solve \\
  PC type  & \texttt{lu} & LU factorization \\
  Damping  & Rayleigh & $\xi = 5\%$ at $(T_1, T_3) = (0.5,\, 0.1)\,$s \\
  Lateral amplitude & 0.005\,MN & Per node per storey \\
  Excitation period & 0.5\,s & Near first-mode period \\
\bottomrule
\end{tabular}
\end{table}

A technical detail of the implementation is that the TS-specific
options are set \emph{programmatically after} \texttt{solver.setup()},
because \texttt{TSSetFromOptions()} (called internally by
\texttt{setup()}) would otherwise overwrite them from the PETSc options
database:
\begin{lstlisting}
solver.setup();
{
    TS ts = solver.get_ts();
    TSAlpha2SetRadius(ts, 0.9);
    TSAdapt adapt;
    TSGetAdapt(ts, &adapt);
    TSAdaptSetType(adapt, TSADAPTNONE);
    TSSetMaxSNESFailures(ts, 5);
}
\end{lstlisting}


% ══════════════════════════════════════════════════════════════════════════
\section{VTK Displacement Reconstruction for Corotational Beams}
\label{sec:vtk-displacement-fix}

\subsection{The Problem: Disconnected Elements in Warp-by-Vector}

When the VTK multi-block exporter was first used with corotational beam
elements, ParaView's \emph{Warp by Vector} filter showed every beam
element as a disconnected fragment: each element appeared to deform
independently, with gaps at every shared node.

The root cause lies in how displacement was reconstructed for the VTK
carrier.  The original code in \texttt{StructuralVTMExporter} computed
displacement at a sampling point~$\xi$ as
\begin{equation}\label{eq:old-displacement}
  \mathbf{u}_\mathrm{vis}(\xi)
  = \mathbf{R}_n^{\mathsf{T}} \,
    \bigl[\mathbf{u}_\mathrm{center,loc}(\xi)
    + \bm{\theta}_\mathrm{loc}(\xi) \times \mathbf{d}_\mathrm{offset}\bigr],
\end{equation}
where $\mathbf{u}_\mathrm{center,loc}$ is the centerline translation in the
\emph{local corotated frame} and $\bm{\theta}_\mathrm{loc}$ is the local
rotation vector.

For \texttt{beam::SmallRotation}, $\mathbf{u}_\mathrm{center,loc}$ contains
the full translational field (small-rotation $\mathbf{T}$ simply rotates
global DOFs into the local frame).  The formula
$\mathbf{R}^{\mathsf{T}} \cdot \mathbf{u}_\mathrm{loc}$ correctly recovers the
global displacement.

For \texttt{beam::Corotational}, however,
$\mathbf{u}_\mathrm{center,loc}$ contains only the \emph{deformational}
part computed from \cref{eq:deformational-dofs}: node~1 translations are
$(0,0,0)$ and node~2 translations are $(\Delta l, 0, 0)$.  The
rigid-body displacement has been deliberately stripped.
Formula~\eqref{eq:old-displacement} therefore returns only the
deformational contribution, which is typically several orders of
magnitude smaller than the actual nodal displacement.

Since each element independently produces a tiny deformational
displacement relative to its own corotated frame, the Warp filter offsets
every element by a different vector --- and the mesh \emph{appears
disconnected}.

% ------------------------------------------------------------------
\subsection{The Fix: Global DOF Interpolation}
\label{sec:global-dof-interpolation}

The fix introduces a new method in \texttt{BeamElement}:
\begin{lstlisting}
Eigen::Vector3d sample_displacement_global(double xi, Vec u_petsc) const
{
    Eigen::VectorXd u_e = extract_element_dofs(u_petsc);
    const double n1 = N1(xi), n2 = N2(xi);
    Eigen::Vector3d u_glob = Eigen::Vector3d::Zero();
    u_glob[0] = n1 * u_e[0] + n2 * u_e[6];
    u_glob[1] = n1 * u_e[1] + n2 * u_e[7];
    u_glob[2] = n1 * u_e[2] + n2 * u_e[8];
    return u_glob;
}
\end{lstlisting}

This method extracts the \emph{global} element DOFs directly from the
PETSc solution vector and interpolates them using the standard linear
shape functions $N_1(\xi), N_2(\xi)$.  Because these are the actual
nodal displacements (not the corotation-filtered deformational DOFs),
the resulting field is continuous across element boundaries.

At an off-axis sampling point with section offset
$\mathbf{d}_\mathrm{off} = (0, y, z)$, the corrected formula becomes
\begin{equation}\label{eq:new-displacement}
  \mathbf{u}_\mathrm{vis}(\xi, y, z)
  = \underbrace{N_1(\xi)\,\mathbf{u}_1 + N_2(\xi)\,\mathbf{u}_2
    \vphantom{\big|}}_{\text{rigid + deformational (global)}}
  + \;
  \underbrace{\mathbf{R}_n^{\mathsf{T}} \bigl[
      \bm{\theta}_\mathrm{loc}(\xi) \times \mathbf{d}_\mathrm{off}
    \bigr]}_{\text{section rotation contribution}},
\end{equation}
where $\mathbf{u}_1, \mathbf{u}_2$ are the global nodal displacement
vectors.  The first term captures the total (rigid + deformational)
centerline motion; the second term adds the rotation-dependent offset
that varies across the cross-section.

% ------------------------------------------------------------------
\subsection{Affected Reconstruction Sites}

The \texttt{StructuralVTMExporter} produces five VTK blocks for each
structural model.  \Cref{tab:vtk-fix-sites} lists every displacement
reconstruction site that was corrected.

\begin{table}[htbp]
\centering
\caption{Displacement reconstruction sites corrected for corotational
  compatibility.  ``Single'' refers to the homogeneous-element exporter;
  ``mixed'' to the \texttt{StructuralElement}-based exporter.}
\label{tab:vtk-fix-sites}
\begin{tabular}{llll}
\toprule
  \textbf{Block} & \textbf{Exporter} & \textbf{Method} & \textbf{Description} \\
\midrule
  axis            & single & \texttt{append\_station}
    & Centerline displacement via \texttt{sample\_displacement\_global} \\
  axis            & mixed  & \texttt{append\_station}
    & Same---dispatched through \texttt{dispatch\_supported\_structural} \\
  surface         & single & \texttt{reconstruct\_beam\_point}
    & Off-axis displacement with section offset \\
  surface         & mixed  & \texttt{reconstruct\_beam\_point}
    & Same---via \texttt{detail::reconstruct\_beam\_point} \\
  sections        & single & \texttt{reconstruct\_beam\_point}
    & Section profile points at each station \\
  sections        & mixed  & \texttt{reconstruct\_beam\_point}
    & Same---dispatched \\
  fibers          & single & inline $\mathbf{R}^{\mathsf{T}} \cdot \mathbf{u}_\mathrm{loc}$
    & Per-fiber displacement at Gauss points \\
  fibers          & mixed  & inline $\mathbf{R}^{\mathsf{T}} \cdot \mathbf{u}_\mathrm{loc}$
    & Same---dispatched \\
\bottomrule
\end{tabular}
\end{table}

The fix for the surface and sections blocks was centralised in the
\texttt{reconstruct\_beam\_point} helper function by adding an optional
\texttt{Vec} parameter.  When a non-null PETSc vector is provided and
the element supports \texttt{sample\_displacement\_global}, the corrected
formula~\eqref{eq:new-displacement} is used:
\begin{lstlisting}
if constexpr (requires { element.sample_displacement_global(xi, u_petsc); }) {
    if (u_petsc) {
        out.displacement = element.sample_displacement_global(xi, u_petsc)
            + R.transpose() * theta.cross(offset_local);
    } else {
        out.displacement = R.transpose() * u_local;
    }
} else {
    out.displacement = R.transpose() * u_local;
}
\end{lstlisting}

The \texttt{if constexpr} / \texttt{requires} guard ensures zero overhead
for formulations where the original reconstruction is correct (e.g.\
shells, small-rotation beams).  The runtime null-check on the PETSc
vector allows unit tests---which use a mock model with
\texttt{state\_vector() = nullptr}---to fall back transparently.

% ------------------------------------------------------------------
\subsection{Type Dispatch for Corotational Elements}

The mixed exporter's \texttt{dispatch\_supported\_structural} function,
which performs runtime type recovery from \texttt{StructuralElement},
was extended with the corotational beam type:
\begin{lstlisting}
// StructuralVTMExporter.hh — dispatch_supported_structural
else if (auto* p = wrapped.template as<
    BeamElement<TimoshenkoBeam3D, 3, beam::Corotational>>()) {
    fn(*p);
}
\end{lstlisting}

Without this entry the exporter threw
\texttt{std::runtime\_error} because the element type name (including
the full \texttt{beam::Corotational} template argument) did not match
any registered concrete type.


% ══════════════════════════════════════════════════════════════════════════
\section{Kinematic Incompatibility: Linear Shell vs.\ Corotational Beam}
\label{sec:kinematic-incompatibility}

\subsection{Diagnosis: Numerical Divergence in the Mixed Model}

Despite correct convergence of SNES at each time step, the simulation
exhibited growing displacements during the gravity ramp, with the
maximum lateral displacement exceeding several metres within a fraction
of a second---far beyond physically reasonable deformation for the
applied load level.

Investigation revealed that the source of instability is not the
corotational beam formulation itself, but rather the
\textbf{kinematic incompatibility between the geometrically nonlinear
beams and the geometrically linear shells} at shared nodes.

% ------------------------------------------------------------------
\subsection{Root Cause Analysis}

\Cref{tab:kinematic-comparison} contrasts the relevant kinematic
assumptions.

\begin{table}[htbp]
\centering
\caption{Kinematic behaviour comparison at shared beam-shell nodes.}
\label{tab:kinematic-comparison}
\begin{tabular}{lcc}
\toprule
  \textbf{Aspect}
  & \textbf{BeamElement (Corotational)}
  & \textbf{ShellElement (Linear)} \\
\midrule
  Local frame $\mathbf{R}$
    & Updated every Newton iteration
    & Fixed at reference configuration \\[3pt]
  Strain measure
    & Corotational local (small in rotating frame)
    & Infinitesimal ($\bm{\varepsilon} = \mathbf{B}_0 \, \mathbf{R}_0 \, \mathbf{u}_e$) \\[3pt]
  Geometric stiffness $\mathbf{K}_\sigma$
    & Yes
    & No \\[3pt]
  Transformation $\mathbf{T}$
    & $\mathbf{T}(\mathbf{R}_n)$, updated
    & $\mathbf{T}(\mathbf{R}_0)$, fixed \\[3pt]
  $\mathbf{f}_\mathrm{int}$ for large $\mathbf{u}_e$
    & Physically correct
    & Linearly extrapolated \\
\bottomrule
\end{tabular}
\end{table}

The fundamental problem is that when the beams undergo large rotations,
the shared nodes move significantly.  The shell elements connected to
those same nodes see large nodal displacement vectors
$\mathbf{u}_e$, but compute their internal forces as
\begin{equation}\label{eq:shell-force}
  \mathbf{f}_\mathrm{int}^\mathrm{shell}
  = \mathbf{T}_0^{\mathsf{T}}
    \sum_{g} w_g \, \lvert J \rvert_g \;
    \mathbf{B}_g^{\mathsf{T}} \,
    \bm{\sigma}\!\bigl(\underbrace{\mathbf{B}_g \, \mathbf{T}_0 \, \mathbf{u}_e
      }_{\text{linearly extrapolated strain}}\bigr).
\end{equation}

For an elastic shell with stiffness matrix $\mathbf{D}$, this reduces to
\begin{equation}
  \mathbf{f}_\mathrm{int}^\mathrm{shell}
  \;=\;
  \underbrace{\mathbf{T}_0^{\mathsf{T}}
    \Bigl(\sum_g w_g \, \lvert J \rvert_g \;
    \mathbf{B}_g^{\mathsf{T}} \, \mathbf{D} \, \mathbf{B}_g\Bigr)
    \mathbf{T}_0}_{\displaystyle \mathbf{K}_0}
  \;\mathbf{u}_e.
\end{equation}

The internal force grows linearly with $\mathbf{u}_e$, without any
geometric correction.  Specifically:
\begin{itemize}
  \item Rigid-body rotations of the shell generate \emph{spurious}
        membrane and bending strains because the reference-frame $B$-matrix
        does not distinguish rotation from deformation.
  \item There is no geometric stiffness contribution from the shell, so
        any destabilising $P$-$\Delta$ effect at the beam-shell junction is
        not captured.
  \item The equilibrium at shared nodes is therefore inconsistent: the
        beam provides geometrically correct forces while the shell provides
        linearly extrapolated forces.
\end{itemize}

% ------------------------------------------------------------------
\subsection{Physical Interpretation}

Consider a storey slab that rigid-body translates laterally by
$u_x = 50$\,mm due to inter-storey drift.  In the corotational beam
formulation, this translation is absorbed by the frame update and
produces no spurious strain.  In the linear shell formulation, however,
the same 50\,mm lateral displacement of all four shell nodes generates a
fictitious membrane strain of order
\begin{equation}
  \varepsilon_\mathrm{spurious}
  \;\approx\; \frac{u_x}{L_\mathrm{panel}}
  \;=\; \frac{0.05}{6.0}
  \;\approx\; 8.3 \times 10^{-3},
\end{equation}
which for a Young's modulus $E \approx 28{,}000$\,MPa yields a spurious
membrane stress of $\sigma_\mathrm{spurious} \approx 233$\,MPa---clearly
nonphysical for an elastic slab subjected only to lateral drift.

The resulting spurious forces from the shell ``push back'' against the
beam displacements, creating an artificial stiffening effect initially.
As the beam elements rotate further and the shell forces grow linearly,
the force imbalance can trigger oscillatory behaviour or outright
divergence in the time integrator.

% ------------------------------------------------------------------
\subsection{Formal Statement of the Incompatibility}

More formally, the corotational beam's internal force satisfies the
\emph{objectivity} requirement:
\begin{equation}\label{eq:objectivity}
  \mathbf{f}_\mathrm{int}^\mathrm{beam}(\mathbf{Q}\,\mathbf{u})
  = \mathbf{Q} \, \mathbf{f}_\mathrm{int}^\mathrm{beam}(\mathbf{u})
\end{equation}
for any superposed rigid-body rotation $\mathbf{Q}$, because the
corotational filter removes $\mathbf{Q}$ before computing local strains.

The linear shell, in contrast, does not satisfy objectivity:
\begin{equation}
  \mathbf{f}_\mathrm{int}^\mathrm{shell}(\mathbf{Q}\,\mathbf{u})
  = \mathbf{K}_0 \, \mathbf{Q} \, \mathbf{u}
  \;\neq\;
  \mathbf{Q} \, \mathbf{K}_0 \, \mathbf{u}
  \quad (\text{in general}).
\end{equation}

At shared nodes, the global equilibrium requires
\begin{equation}
  \mathbf{f}_\mathrm{int}^\mathrm{beam}
  + \mathbf{f}_\mathrm{int}^\mathrm{shell}
  + \mathbf{f}_\mathrm{inertia}
  + \mathbf{f}_\mathrm{damping}
  = \mathbf{f}_\mathrm{ext},
\end{equation}
but the lack of objectivity in the shell term introduces
\emph{configuration-dependent} artificial forces that pollute the
equilibrium.  The magnitude of this pollution grows with the deviation
from the reference configuration.


% ══════════════════════════════════════════════════════════════════════════
\section{Resolution Strategies}
\label{sec:resolution-strategies}

Three paths can resolve the kinematic incompatibility:

\begin{description}
  \item[Option A: Corotational shell formulation.]
    Implement \texttt{beam::Corotational}-style frame update for
    \texttt{ShellElement}: at each Newton iteration, compute the
    mean rotation of the element from the current nodal positions and
    project out the rigid-body component.  This is the theoretically
    clean solution and makes the shell objective under superposed
    rigid-body motions.
    \emph{Estimated complexity:} medium---the corotational machinery
    is already available from the beam policy.

  \item[Option B: Keep shells linear, limit deformations.]
    Restrict the problem to load levels where inter-storey drift is
    below $\approx\!1\%$, so that the shell's linear kinematic error
    remains negligible.  This is the pragmatic path already used in the
    previous small-rotation static analysis.

  \item[Option C: Replace slabs with rigid diaphragm constraints.]
    Many earthquake engineering codes allow modelling floor slabs as
    rigid diaphragms---each floor level shares a single translational
    DOF in the horizontal plane plus a single rotation about the
    vertical axis.  This eliminates the shell elements entirely and
    removes the kinematic incompatibility at the cost of losing in-plane
    slab flexibility.
\end{description}

Option~A is the long-term architectural target.
Chapter~\ref{ch:regular-5story-rc-nonlinear-building} documented the
small-rotation version (effectively Option~B).

For immediate use, the lateral load amplitude was reduced to
$0.005$\,MN/node (5\,kN) and the gravity was ramped slowly to
avoid exciting the stiff shell response, keeping drifts within the
linear-shell validity range.

% ══════════════════════════════════════════════════════════════════════════
\section{Test Regression Summary}
\label{sec:corot-test-regression}

\Cref{tab:corot-test-results} summarises the test status after all
modifications described in this chapter.

\begin{table}[htbp]
\centering
\caption{CTest results after corotational VTK fixes.
  Pre-existing failures are due to missing mesh files, not code defects.}
\label{tab:corot-test-results}
\begin{tabular}{lr}
\toprule
  \textbf{Metric} & \textbf{Value} \\
\midrule
  Total tests         & 47 \\
  Passed              & 45 \\
  Failed              & 2 (pre-existing) \\
  New regressions     & 0 \\
\midrule
  \texttt{structural\_vtm\_export}       & \textbf{pass} \\
  \texttt{structural\_mixed\_vtm\_export} & \textbf{pass} \\
\bottomrule
\end{tabular}
\end{table}

The two pre-existing failures (\texttt{cantilever\_benchmark} and
\texttt{snes\_gmsh\_pipeline}) are due to missing \texttt{.msh} input
files that were never committed to the repository.

% ══════════════════════════════════════════════════════════════════════════
\section{Files Modified}
\label{sec:corot-files-modified}

\begin{description}
  \item[\texttt{src/elements/BeamElement.hh}.]
    Added \texttt{sample\_displacement\_global(xi, Vec)} method
    for direct global-DOF interpolation.

  \item[\texttt{src/post-processing/VTK/StructuralVTMExporter.hh}.]
    (i) Extended \texttt{dispatch\_supported\_structural} with
    \texttt{BeamElement<\ldots, beam::Corotational>}.
    (ii) Modified \texttt{reconstruct\_beam\_point} to accept an
    optional \texttt{Vec} and apply \cref{eq:new-displacement}.
    (iii) Modified axis and fiber blocks in both single-element and
    mixed exporters with \texttt{if constexpr} / null-guard idiom.

  \item[\texttt{tests/test\_structural\_mixed\_vtm\_export.cpp}.]
    Added \texttt{is\_beam\_element} catch-all branch for corotational
    beam type to return 12-DOF zero vector instead of 24-DOF shell
    vector.

  \item[\texttt{main.cpp}.]
    (i) Changed element type from \texttt{beam::SmallRotation} to
    \texttt{beam::Corotational}.
    (ii) Added gravity/lateral ramp functions.
    (iii) Reduced lateral amplitude to 0.005\,MN/node.
    (iv) Moved TS options to programmatic API after
    \texttt{solver.setup()}.
\end{description}

% ══════════════════════════════════════════════════════════════════════════
\section{Summary}

This chapter documented three interrelated developments:
\begin{enumerate}
  \item The corotational beam kinematic policy
        (\texttt{beam::Corotational}) enables geometrically nonlinear
        analysis of the frame while retaining the small-strain
        $B$-matrix in a rotating local frame.

  \item The VTK displacement reconstruction was incorrect for
        corotational elements because the deformational-only local DOFs
        were being used where total global displacements were required.
        The fix introduces \texttt{sample\_displacement\_global} and
        applies it across all eight reconstruction sites with
        compile-time guards and runtime null checks.

  \item The mixing of geometrically linear Mindlin shells with
        corotational beams creates a fundamental kinematic
        incompatibility: rigid-body rotations at shared nodes induce
        spurious strains in the linear shells, violating
        objectivity~\eqref{eq:objectivity} and destabilising the global
        equilibrium.  The long-term resolution is a corotational shell
        formulation.
\end{enumerate}
